\begin{table}[htbp]
\centering
\footnotesize\setlength{\tabcolsep}{4pt}
\caption{Comparación de modelos anidados (misma muestra y mismos pesos)}\label{tab:comparacion}
\begin{tabularx}{\linewidth}{>{\raggedright\arraybackslash}X r r r r r r r r}
\toprule
Modelo & $k$ & $\bar R^2$ & AIC & BIC & \multicolumn{2}{c}{$F$ clásico} & \multicolumn{2}{c}{Wald robusto} \\ \cmidrule(lr){6-7}\cmidrule(lr){8-9} & & & & & $F$ & $p$ & $F$ & $p$ \\
\midrule
Nulo (región) & 3 & \num{0.154} & \num{25083} & \num{25107} & --- & --- & --- & --- \\
Máximo (tras la poda) & 24 & \num{0.544} & \num{23436} & \num{23584} & \num{110.70} & $<\num{0.001}$ & \num{106.52} & $<\num{0.001}$ \\
Seleccionado (efectos principales) & 13 & \num{0.537} & \num{23462} & \num{23545} & \num{4.35} & $<\num{0.001}$ & \num{2.13} & \num{0.036} \\
Final (con interacciones) & 19 & \num{0.543} & \num{23433} & \num{23551} & \num{6.85} & $<\num{0.001}$ & \num{5.71} & $<\num{0.001}$ \\
\bottomrule
\end{tabularx}
\par\smallskip{\footnotesize\raggedright Cada fila se compara con la anterior (bloque de términos que se añade o se retira). El $F$ clásico supone errores independientes y homocedásticos; el Wald robusto usa la covarianza por conglomerados de estado, la misma que guía la selección.\par}
\end{table}
